% Pursuing Stacks (cotes 134-2 to 134-8): archive cote 134-5, site page 56 (PDF page 57), typescript page 243.
% typescript. Draft compared: ps4.tex:2830-2896.
% AI-transcribed by Claude Opus 5.5 (claude-opus-5-5, adaptive thinking, effort high), Message Batches API, batch ps2 (msgbatch_01NVwr1Yuo8qJf5rykeGBXPK), 2026-11-12; reviewed against both supplied scan halves; see CHECK_DECISIONS.tsv.
\scanpage{134-5/56}
\mspage{243}
corollary,  when $M = (\mathrm{Cat})$ ,  $M_{as} = (\mathrm{Cat})_{as}$ ).  Now this follows at once
from \inserted[typed]{strong} saturation of $\underline{W}$ ,  which (as we announced earlier) follows from
Loc 1)  to  Loc 3)  (and will be pro\xout{\unsure{s}}ved in \xout{the} part \hand{V} of the notes).
The reader who fears a vicious circle may till then restrict \xout{eventual}
use of the theorem to the case when we assume beforehand that $\over{\illegible{1}}{e}_M$ and
$e_{A^\wedge}$ are aspheric .

It is now clear that the six conditions of prop. 4  are equiva-%
lent,  and they \xout{imply} \inserted[typed]{are} of course \xout{\unsure{(iii)}} implied by (iii).  Conversely,
they imply (iii),  as follows from the fact that in the canonical diagram
(commutative up to canonical isomorphism)
% Hand-drawn arrows are marked blue; typed objects remain plain.
\diagram{\begin{tikzcd}[column sep=tiny]
\mathrm{Hot}_M \ar[rr,color=blue] \ar[dr,color=blue] \& \& \mathrm{Hot}_A \ar[dl,color=blue] \\
\& (\mathrm{Hot}) \&
\end{tikzcd}}
,
the two downwards arrows are equivalences,  qed.

\xout{The theorem}
The theorem above seems to me to be exactly the ``something very
simple-minded surely'' which I was feeling to get b\unsure{ur}ningly clos\unsure{e}, by
the end of march,  ne\over[typed]{q}{a}rly three months ago (p.89);  at least, \inserted[typed]{to be ``just it'' ,} as far as
the case of \hand{\uline{\textcolor{black}{test-functors}}} is concerned.  We may equally view th\over{\unsure{i}}{i}s theo-%
rem as giving the precise relationship between the notion of a weak test
functor or a \hand{\uline{\textcolor{black}{test functor}}}, \hand{(}the latest version of which (in the \unsure{c}ontext
of $\underline{W}$-notions)  appears in section 65  (def.5 and 6,  p. 175 and 176)\hand{)},
and the notion of \hand{\uline{\textcolor{black}{aspheric functors}}},  more precisely of $M_{as}$-aspheric
functors,  introduced lately (p. 233).  This now is the moment surely
to check if the terminology of test functors and weak ones  introduced
earlier,  before the relevant notion of asphericity structures was at
hand,  is really satisfactory indeed,  and if needed,  adjust it slightly.

So let again
\[ i\colon A \hand{\longrightarrow} M \]
be a functor,  with A small,  and $(M,M_{as})$ an \inserted{\deleted{aspheric}} asphericity structure.  We
don't  assume beforehand,  neither that A is a test-category or the like,
nor that $(M,M_{as})$ be modelizing.  We now paraphrase def.5  (p.175) of
weak test functors as follows:

\hand{\uline{\textcolor{black}{Definition}}} \deleted{\illegible{1}}1:  The functor i above is called a \hand{\uline{\textcolor{black}{pse\unsure{u|d}do-test}}} \hand{\uline{\textcolor{black}{functor}}} ,
if it satisfies the following three conditions:

a)  the corresponding functor  $i^{*}\colon M \hand{\to} A^\wedge$ is ``model-preserving'',  by

\margin{56}
